% !TEX program = xelatex
\documentclass[11pt,a4paper]{article}
\usepackage{algo-preamble}

\title{\textbf{L17 \textperiodcentered{} Δυναμικός Προγραμματισμός IV}\\[2pt]
\large Ανεξάρτητο Σύνολο σε Δέντρα \& Bellman-Ford}
\author{Algorithms Class Hub}
\date{\small Αλγόριθμοι και Πολυπλοκότητα (K17) \textperiodcentered{} ΕΚΠΑ}

\begin{document}
\maketitle

\begin{center}\itshape\small
Δυναμικός προγραμματισμός σε γραφήματα: ανεξάρτητο σύνολο μέγιστου βάρους σε
δέντρα, και ο αλγόριθμος Bellman-Ford για συντομότερες διαδρομές με αρνητικά
βάρη.
\end{center}

\section{Τι θα δούμε}

Η τελευταία διάλεξη του μαθήματος εφαρμόζει τον δυναμικό προγραμματισμό σε
\textbf{γραφήματα}:
\begin{enumerate}
  \item \textbf{Ανεξάρτητο σύνολο σε δέντρα} --- DP όπου τα υποπροβλήματα είναι
  \textbf{υποδέντρα}.
  \item \textbf{Bellman-Ford} --- συντομότερες διαδρομές με \textbf{αρνητικά}
  βάρη, εκεί που ο Dijkstra (από το μάθημα των γραφημάτων) αποτυγχάνει. Κλείνει
  έτσι μια υπόσχεση που είχαμε αφήσει ανοιχτή.
\end{enumerate}

\section{Ανεξάρτητο Σύνολο σε Δέντρα --- «το πάρτι»}

\subsection{Το πρόβλημα}

Οργανώνουμε ένα πάρτι για μια εταιρεία με \textbf{αυστηρή ιεραρχική δομή}. Κάθε
υπάλληλος $v$ έχει μια θετική «προσφορά» $\chi(v)$ (πόσο κέφι φέρνει). Ένας
περιορισμός: \textbf{δεν} καλούμε ταυτόχρονα έναν υπάλληλο και τον \textbf{άμεσο
προϊστάμενό} του.

\textbf{Είσοδος:} ένα \textbf{δέντρο} που περιγράφει την ιεραρχία, με προσφορά σε
κάθε κορυφή.

\textbf{Στόχος:} λίστα καλεσμένων με το \textbf{μέγιστο άθροισμα προσφορών}, που
να σέβεται τον περιορισμό.

\begin{intuition}
Σε ορολογία γραφημάτων: ζητάμε ένα \textbf{ανεξάρτητο σύνολο} (σύνολο κορυφών
χωρίς ακμή μεταξύ τους) \textbf{μέγιστου βάρους} σε ένα δέντρο. Στα γενικά
γραφήματα αυτό το πρόβλημα είναι NP-δύσκολο --- αλλά σε \textbf{δέντρα} ο
δυναμικός προγραμματισμός το λύνει εύκολα.
\end{intuition}

\subsection{Δύο τιμές ανά κορυφή}

Το κλειδί: για κάθε κορυφή $v$, τα υποπροβλήματα αφορούν το \textbf{υποδέντρο}
της $v$. Αλλά μία τιμή δεν αρκεί --- η απάντηση εξαρτάται από το \textbf{αν η
ίδια η $v$ είναι μέσα}. Άρα κρατάμε \textbf{δύο}:

\begin{keypoint}
\begin{itemize}
  \item $A[v] = $ μέγιστη προσφορά ενός έγκυρου υποσυνόλου του \textbf{υποδέντρου}
  της $v$ (η $v$ επιτρέπεται ή όχι).
  \item $B[v] = $ μέγιστη προσφορά ενός έγκυρου υποσυνόλου του υποδέντρου της $v$
  που \textbf{δεν περιέχει} την $v$.
\end{itemize}
\end{keypoint}

Έστω $v$ με άμεσους υφισταμένους (παιδιά) $v_1, v_2, \dots, v_k$.

\textbf{Για το $B[v]$:} η $v$ \textbf{δεν} μπαίνει --- άρα δεν μπλοκάρει
κανέναν. Κάθε υφιστάμενος μπορεί ελεύθερα να μπει ή όχι· για κάθε υποδέντρο
παίρνουμε το βέλτιστο $A[v_i]$:
\[
B[v] = \sum_{i=1}^{k} A[v_i]
\]

\textbf{Για το $A[v]$:} δύο σενάρια, παίρνουμε το καλύτερο. Η $v$ \textbf{δεν}
μπαίνει: αυτό είναι ακριβώς το $B[v]$. Η $v$ \textbf{μπαίνει}: τότε
\textbf{κανένα} από τα παιδιά $v_1, \dots, v_k$ δεν επιτρέπεται, με συνεισφορά
$\chi(v) + \sum_i B[v_i]$.
\[
A[v] = \max\Bigl\{\, B[v],\ \ \chi(v) + \sum_{i=1}^{k} B[v_i] \,\Bigr\}
\]

\begin{notebox}
\textbf{Παράδειγμα.} Σε ένα μικρό δέντρο με ρίζα $r$ ($\chi = 8$) και παιδιά $a$
($\chi = 10$) και $b$ ($\chi = 20$), όπου το $a$ έχει παιδιά $c$ ($\chi = 15$)
και $d$ ($\chi = 4$), η βέλτιστη λίστα είναι $\{b, c, d\}$ με άθροισμα
$20 + 15 + 4 = 39$. Πράγματι $A[r] = \max\{B[r] = 39,\ \chi(r) + B[a] + B[b] =
8 + 19 + 9 = 36\} = 39$.
\end{notebox}

\subsection{Ο αλγόριθμος}

Υπολογίζουμε τα $A[v], B[v]$ \textbf{από τα φύλλα προς τη ρίζα} (κάθε κορυφή
χρειάζεται πρώτα τα παιδιά της --- μια αντίστροφη τοπολογική διάταξη του
δέντρου). Για \textbf{φύλλο} $\ell$: $B[\ell] = 0$ και $A[\ell] = \chi(\ell)$. Η
απάντηση είναι το $A[\text{ρίζα}]$.

Κάθε κορυφή και κάθε ακμή επεξεργάζεται $O(1)$ φορές $\to$ \textbf{$O(n)$}
συνολικά.

\section{Συντομότερες Διαδρομές με Αρνητικά Βάρη}

\subsection{Γιατί ξαναπιάνουμε το πρόβλημα}

Στο μάθημα των γραφημάτων ο Dijkstra έλυσε τις συντομότερες διαδρομές --- αλλά με
τη ρητή \textbf{απαίτηση όλα τα βάρη να είναι θετικά}. Τώρα επιτρέπουμε
\textbf{αρνητικά} βάρη $\ell_e$ (π.χ.\ δρομολόγηση όπου κάποιες ακμές δίνουν
«κέρδος»).

\begin{warningbox}
\textbf{Ο Dijkstra μπορεί να αποτύχει.} Μόλις «οριστικοποιήσει» μια κορυφή, δεν
την ξανακοιτάζει. Μια αρνητική ακμή που οδηγεί σε \textbf{ήδη οριστικοποιημένη}
κορυφή θα την βελτίωνε --- αλλά είναι πια αργά. (Παράδειγμα: με ακμές
$s \to t$ βάρους 2, $s \to u$ βάρους 1 και $t \to u$ βάρους $-5$, ο Dijkstra
οριστικοποιεί το $u$ στο 1 μέσω $s \to u$ και χάνει το
$s \to t \to u = 2 + (-5) = -3$.)
\end{warningbox}

\subsection{Πρόσθεσε σταθερά; Όχι.}

Δελεαστική ιδέα: «πρόσθεσε μια μεγάλη σταθερά σε \textbf{όλα} τα βάρη ώστε να
γίνουν θετικά, και τρέξε Dijkstra». \textbf{Δεν δουλεύει:} μια διαδρομή με
πολλές ακμές κερδίζει τη σταθερά \textbf{πολλές φορές}, ενώ μια διαδρομή με λίγες
ακμές λίγες --- η σύγκριση αλλοιώνεται. Ο Dijkstra και πάλι αποτυγχάνει.

\subsection{Αρνητικοί κύκλοι}

\begin{keypoint}
\textbf{Αρνητικός κύκλος:} ένας κύκλος του οποίου το άθροισμα βαρών είναι
\textbf{αρνητικό}.

Αν υπάρχει διαδρομή από την $s$ στην $t$ που περνά μέσα από αρνητικό κύκλο,
\textbf{δεν υπάρχει} συντομότερη s--t διαδρομή --- γυρνώντας στον κύκλο μειώνεις
το κόστος επ' άπειρον. Διαφορετικά, η συντομότερη διαδρομή \textbf{υπάρχει} και
είναι \textbf{απλή} (χωρίς επαναλαμβανόμενες κορυφές).
\end{keypoint}

\section{Ο αλγόριθμος Bellman-Ford}

\subsection{Η αναδρομή --- παραμετροποίηση με το πλήθος ακμών}

Σταθεροποιούμε έναν \textbf{προορισμό} $t$. Το έξυπνο υποπρόβλημα:

\begin{keypoint}
$\text{OPT}(i, v) = $ το μήκος του συντομότερου $v \to t$ μονοπατιού που
χρησιμοποιεί \textbf{το πολύ $i$ ακμές}.
\end{keypoint}

Έστω $P$ το βέλτιστο τέτοιο μονοπάτι. Δύο περιπτώσεις:
\begin{itemize}
  \item Το $P$ χρησιμοποιεί \textbf{το πολύ $i-1$} ακμές --- τότε
  $\text{OPT}(i, v) = \text{OPT}(i-1, v)$.
  \item Το $P$ χρησιμοποιεί \textbf{ακριβώς $i$} ακμές. Έστω $(v, w)$ η πρώτη
  του ακμή· τότε το υπόλοιπο είναι το συντομότερο $w \to t$ μονοπάτι με $i-1$
  ακμές. Κόστος: $\ell_{vw} + \text{OPT}(i-1, w)$.
\end{itemize}

\begin{keypoint}
\[
\text{OPT}(i, v) =
\begin{cases}
0 & i = 0,\ v = t \\
\infty & i = 0,\ v \ne t \\[4pt]
\min\Bigl\{\, \text{OPT}(i-1, v),\ \ \min_{(v,w) \in A}\bigl(\ell_{vw} +
\text{OPT}(i-1, w)\bigr) \,\Bigr\} & i > 0
\end{cases}
\]
\end{keypoint}

\begin{intuition}
\textbf{Πόσο μεγάλο πρέπει να γίνει το $i$;} Αν δεν υπάρχουν αρνητικοί κύκλοι, η
συντομότερη διαδρομή είναι \textbf{απλή} --- δεν επαναλαμβάνει κορυφή --- άρα έχει
\textbf{το πολύ $n-1$ ακμές}. Συνεπώς το $M[n-1, v]$ δίνει τη συντομότερη
$v \to t$ απόσταση για \textbf{κάθε} κορυφή $v$.
\end{intuition}

\subsection{Ο πίνακας}

\begin{codebox}
\begin{verbatim}
ΣυντομότεροΜονοπάτι(G, t):
  Για κάθε v ∈ V:  M[0, v] <- ∞
  M[0, t] <- 0
  Για i = 1 έως n-1:
    Για κάθε v ∈ V:
      M[i, v] <- M[i-1, v]
      Για κάθε ακμή (v, w) ∈ A:
        M[i, v] <- min( M[i, v],  l(v,w) + M[i-1, w] )
  επίστρεψε M[n-1, .]
\end{verbatim}
\end{codebox}

\textbf{Πολυπλοκότητα:} $\Theta(mn)$ χρόνος, $\Theta(n^2)$ χώρος.

\subsection{Βελτιώσεις}

\begin{keypoint}
\begin{itemize}
  \item \textbf{Γραμμικός χώρος.} Κράτα έναν \textbf{μονοδιάστατο} πίνακα
  $M[v] = $ το μήκος του καλύτερου $v \to t$ μονοπατιού που βρέθηκε ως τώρα.
  Μετά από $i$ γύρους, το $M[v]$ δεν είναι μεγαλύτερο από το συντομότερο
  μονοπάτι με $\le i$ ακμές. Χώρος: $O(m + n)$.
  \item \textbf{Παράλειψη.} Μην εξετάζεις ακμή $(v, w)$ εκτός αν το $M[w]$
  \textbf{άλλαξε} στον προηγούμενο γύρο --- αλλιώς δεν θα δώσει τίποτα νέο.
  \item \textbf{Πρόωρος τερματισμός.} Αν σε έναν ολόκληρο γύρο \textbf{καμία}
  τιμή δεν άλλαξε, σταμάτα --- έχεις συγκλίνει.
\end{itemize}
\end{keypoint}

Με αυτές: $O(m + n)$ χώρος, $O(mn)$ χρόνος στη χειρότερη περίπτωση --- αλλά
συχνά πολύ ταχύτερα.

\begin{notebox}
\textbf{Παρατηρήσεις-ερωτήματα.} Ο αλγόριθμος, δοσμένου ενός $t$, υπολογίζει το
συντομότερο $v \to t$ μονοπάτι για \textbf{κάθε} $v$. Μπορεί να προσαρμοστεί για
$s \to v$ (αντιστρέφοντας τις ακμές), για \textbf{DAG} (όπου αρκεί ένα πέρασμα
κατά τοπολογική διάταξη), και --- αλλάζοντας \texttt{min} σε \texttt{max} ---
για τα \textbf{μακρύτερα} μονοπάτια.
\end{notebox}

\begin{recapbox}
\textbf{Τι κρατάμε από αυτή τη διάλεξη}
\begin{enumerate}
  \item \textbf{Ανεξάρτητο σύνολο σε δέντρα} --- DP με υποπροβλήματα τα
  υποδέντρα. Δύο τιμές ανά κορυφή: $A[v]$ (η $v$ ελεύθερη) και $B[v]$ (η $v$
  έξω). $B[v] = \sum A[v_i]$, $A[v] = \max\{B[v],\ \chi(v) + \sum B[v_i]\}$.
  $O(n)$.
  \item Ο \textbf{Dijkstra αποτυγχάνει} με αρνητικά βάρη· η προσθήκη σταθεράς
  \textbf{δεν} βοηθάει.
  \item \textbf{Αρνητικός κύκλος} $\to$ δεν υπάρχει συντομότερη διαδρομή·
  αλλιώς υπάρχει και είναι \textbf{απλή} ($\le n-1$ ακμές).
  \item \textbf{Bellman-Ford} --- $\text{OPT}(i, v) = $ συντομότερο $v \to t$
  μονοπάτι με $\le i$ ακμές. Αναδρομή: $\min\{\text{OPT}(i-1,v),\
  \min_{(v,w)}(\ell_{vw} + \text{OPT}(i-1,w))\}$. Απάντηση στο $M[n-1, v]$.
  \item $\Theta(mn)$ χρόνος· με μονοδιάστατο πίνακα, παράλειψη αμετάβλητων
  κορυφών και πρόωρο τερματισμό $\to O(m+n)$ χώρος.
\end{enumerate}
\end{recapbox}

\lecturefooter
\end{document}
